\TemplateTodo{Tulis abstrak sesuai bahasa dokumen. Paragraf abstrak memuat konteks, tujuan, metode, data atau objek penelitian, hasil utama, dan kesimpulan. Hindari sitasi dan klaim yang belum dibuktikan.}

\noindent Pertumbuhan pengguna layanan transportasi dan pesan-antar makanan online di Indonesia menghasilkan volume ulasan yang sangat besar di platform seperti Google Play Store. Ulasan tersebut mengandung opini pengguna yang berharga, namun sulit dianalisis secara manual. Penelitian ini bertujuan untuk mengklasifikasikan sentimen ulasan aplikasi GoFood ke dalam kelas positif dan negatif menggunakan metode Na\"ive Bayes Classifier (NBC). Dataset terdiri dari 5.000 ulasan berbahasa Indonesia yang diperoleh melalui teknik web scraping. Tahap preprocessing meliputi case folding, tokenisasi, filtering, dan stemming. Fitur diekstraksi menggunakan TF-IDF (Term Frequency-Inverse Document Frequency). Hasil pengujian dengan 10-fold cross validation menunjukkan bahwa NBC mencapai akurasi sebesar 86,3\%, precision 84,7\%, recall 88,1\%, dan F1-score 86,4\%. Hasil ini menunjukkan bahwa NBC mampu mengklasifikasikan sentimen ulasan dengan baik dan dapat dijadikan acuan untuk pengembangan sistem analisis sentimen yang lebih lanjut.